\begingroup\small\setlength{\tabcolsep}{3pt}
\begin{table}[H]\centering
\caption*{\textbf{表S1　去极值方法敏感性}}
\begin{tabular*}{\textwidth}{@{\extracolsep{\fill}}lrrrrr@{}}\toprule
去极值方案 & \makecell{Rank IC\\均值} & ICIR & \makecell{回归因子收益\\均值（\%/月）} & \makecell{多空年化\\净收益（\%）} & \makecell{Q5算术年化超额\\（百分点）} \\ \midrule
MAD5 & 0.0566 & 0.550 & 0.40 & 13.00 & 5.31 \\
MAD3 & 0.0565 & 0.522 & 0.40 & 12.34 & 4.80 \\
1\%/99\%分位数缩尾 & 0.0563 & 0.563 & 0.38 & 13.32 & 5.55 \\
\bottomrule\end{tabular*}
\par\smallskip\raggedright\footnotesize 全部 A 股，228期；仅改变去极值方法，其余设定及输入S样本不变。ICIR未年化；多空为扣除双腿成本后的CAGR；Q5超额为相对原内部等权毛收益基准的算术年化超额。分位点采用线性插值。差值按未舍入结果计算。
\end{table}
\endgroup
